\begin{center}
    {\LARGE \textbf{2009分析化学}} \\[2ex]
    {\large 时量180分钟 \quad 满分90分}
\end{center}

\vspace{0.5cm}
\noindent\textbf{注意事项:}
\begin{enumerate}[label=\arabic*., parsep=0pt, itemsep=0pt]
    \item 答题前,考生先将自己的姓名、学号、考场号、座位号填写在答题卡上,将条形码准确粘贴在考生信息条形码粘贴区。
    \item 考试结束后,将本试卷与答题纸一并收回。
    \item 回答各个问题时,将答案写在答题纸上,写在本试卷上无效。
\end{enumerate}

\vspace{0.5cm}
\noindent\textbf{一、简答题(28 pts.)}

\begin{enumerate}[label=\arabic*., itemsep=1.5ex]
    \item 在下列物质中：
    
    $\ce{NH4Cl}$（$\mathrm{p}K_\mathrm{b}(\ce{NH3})=4.74$）、
    
    苯酚$\ce{C6H5OH}$（$\mathrm{p}K_\mathrm{a}=9.96$）、
    
    $\ce{Na2CO3}$（$\ce{CH2O3}$的$\mathrm{p}K_{\mathrm{a1}}=6.38$，$\mathrm{p}K_{\mathrm{a2}}=10.25$）、
    
    $\ce{NaAc}$（$\mathrm{p}K_\mathrm{a}(\ce{HAc})=4.74$）、
    
    甲酸$\ce{HCOOH}$（$\mathrm{p}K_\mathrm{a}=3.74$）

    请分别写出能用强碱标准溶液直接滴定的物质以及能用强酸标准溶液直接滴定的物质。
    \item 写出用$\ce{K2Cr2O7}$标准溶液标定$\ce{Na2S2O3}$的两个反应方程式。
    \item 某同学如下配制$0.02\ \mathrm{mol/L}\ \ce{Na2S2O3}$ $500\ \mathrm{mL}$溶液：准确称取$\ce{Na2S2O3\cdot 5H2O}$ $2.482\ \mathrm{g}$，溶于蒸馏水中，加热煮沸，冷却转移至$500\ \mathrm{mL}$容量瓶中，加蒸馏水定容摇匀，保存待用。请指出其错误。
    \item 假定某盐试样为化学纯$\ce{MgSO4\cdot 7H2O}$，称取$0.8000\ \mathrm{g}$试样，将镁沉淀为$\ce{MgNH4PO4}$并灼烧成$\ce{Mg2P2O7}$，得$0.3900\ \mathrm{g}$，试问试样是否符合已知的化学式？原因何在？
    
    已知$M(\ce{MgSO4\cdot 7H2O})=246.5$，$M(\ce{MgSO4})=120.4$，$M(\ce{H2O})=18.02$，$M(\ce{Mg2P2O7})=222.6$
    \item 在滴定分析中，若指示剂的变色点与化学计量点恰好相同，能不能说滴定误差为零，为什么？
    \item 某溶液含有$\ce{M}$和$\ce{N}$两种金属离子，已知$K(\ce{MY}) > K(\ce{NY})$。$\lg K'_{\ce{MY}}$先随溶液的$\mathrm{pH}$增加而增大，为什么？而当$\mathrm{pH}$增加时，$\lg K'_{\ce{MY}}$保持在某一定值（$\ce{N}$在此条件下水解），为什么？
    \item 用佛尔哈德法测定$\ce{Cl-}$时，若不采用加硝基苯等方法，对分析结果有何影响？用法扬司法滴定$\ce{Cl-}$时，若用曙红作指示剂，对分析结果有何影响？
\end{enumerate}

\vspace{0.5cm}
\noindent\textbf{二、计算题(62 pts. total)}

\begin{enumerate}[label=\arabic*., itemsep=1.5ex]
    \item (10 pts.) 以$0.10\ \mathrm{mol/L}\ \ce{NaOH}$溶液滴定$0.10\ \mathrm{mol/L}\ \ce{HCl}$和$0.20\ \mathrm{mol/L}\ \ce{H3BO3}$的混合溶液。
    \begin{enumerate}[label=(\arabic*), noitemsep, topsep=0pt, partopsep=0pt]
        \item 计算化学计量点时溶液的$\mathrm{pH}$；
        \item 若滴定终点$\mathrm{pH}$比化学计量点高$0.5$，计算终点误差。（已知$K_\mathrm{a}(\ce{H3BO3})=5.8\times10^{-10}$）
    \end{enumerate}
    \item (12 pts.) 在$\mathrm{pH}=10.0$氨性缓冲溶液中，以铬黑$\mathrm{T}$为指示剂，用$2.00\times10^{-2}\ \mathrm{mol/L}\ \mathrm{EDTA}$滴定相同浓度的$\ce{Zn^{2+}}$。若终点时游离氨的浓度为$0.20\ \mathrm{mol/L}$，计算终点误差。（已知$\mathrm{pH}=10.0$时，$\lg\alpha_{\mathrm{Y(H)}}=0.45$，$\lg\alpha_{\mathrm{Zn(OH)}}=2.4$，$\mathrm{pZn_{t(\text{铬黑T})}}=12.2$，$\lg K_{\ce{ZnY}}=16.5$；$\ce{Zn^{2+}-NH3}$络合物的$\lg\beta_1 \sim \lg\beta_4$分别为$2.37, 4.81, 7.31, 9.46$）
    \item (10 pts.) 计算$\mathrm{pH}=8.0$时，$\ce{As(V)/As(III)}$电对的条件电位（忽略离子强度的影响），并从计算结果判断以下反应的方向：
    \[
    \ce{H3AsO4 + 2H+ + 3I- <=> HAsO2 + I3- + 2H2O}
    \]
    已知：$E^\ominus_{\ce{H3AsO4/HAsO2}}=0.56\ \mathrm{V}$，$E^\ominus_{\ce{I2/I-}}=0.545\ \mathrm{V}$；$\ce{H3AsO4}$的$K_{\mathrm{a1}}=10^{-2.2}$，$K_{\mathrm{a2}}=10^{-7.0}$，$K_{\mathrm{a3}}=10^{-11.5}$；$\ce{HAsO2}$的$K_\mathrm{a}=10^{-9.22}$。
    \item (8 pts.) 在$\mathrm{pH}=10.0$的$\ce{Ca^{2+}}$、$\ce{Ba^{2+}}$混合试液中加入$\ce{K2CrO4}$，使$[\ce{CrO4^{2-}}]$为$0.010\ \mathrm{mol/L}$。问能否屏蔽$\ce{Ba^{2+}}$而准确滴定$\ce{Ca^{2+}}$？（已知$\mathrm{pH}=10.0$时，$\lg\alpha_{\mathrm{Y(H)}}=0.45$，$\lg K_{\ce{CaY}}=10.69$，$\lg K_{\ce{BaY}}=7.86$，$K_{\mathrm{sp}}(\ce{BaCrO4})=1.2\times10^{-10}$）
    \item (12 pts.) 某方法提出了一新的分析方法，并用此方法测定了一个标准试样，得下列数据$(\%)$（按大小排列）：$40.00, 40.15, 40.16, 40.18, 40.20$。已知该试样的标准值为$40.19\%$（显著水平$0.05$）。
    \begin{enumerate}[label=(\arabic*), noitemsep, topsep=0pt, partopsep=0pt]
        \item 用格鲁布斯（Grubbs）法，检验极端值是否应该舍弃？
        \item 试用$t$检验法对新结果作出评价。
    \end{enumerate}
    附表：
    \begin{center}
    \begin{tabular}{l|ccc}
    $n$ & 4 & 5 & 6 \\
    \hline
    $T_{0.05,n}$ & 1.463 & 1.672 & 1.832 \\
    \end{tabular}
    \begin{tabular}{l|ccc}
    $f$ & 2 & 3 & 4 \\
    \hline
    $t_{0.05,f(\text{双边})}$ & 4.303 & 3.182 & 2.776 \\
    \end{tabular}
    \end{center}
    \item (10分) 某金属离子$\ce{M^{4+}}$与卤离子$\ce{X-}$在弱酸性条件下形成$\ce{MX3}$络合物（不存在其他类型的络合物），该络合物在$372\ \mathrm{nm}$处有最大吸收。今有相同体积的弱酸性试液两份，其中$c(\ce{M^{4+}})$均为$5.0\times10^{-4}\ \mathrm{mol/L}$。
    \begin{enumerate}[label=(\arabic*), noitemsep, topsep=0pt, partopsep=0pt]
        \item 在第一份显色液中$c(\ce{X-})=0.20\ \mathrm{mol/L}$，$\ce{M^{4+}}$显色完全，以$1\ \mathrm{cm}$比色皿，于$372\ \mathrm{nm}$处测得吸光度为$0.748$。
        \item 在第二份显色液中$c(\ce{X-})=2.0\times10^{-3}\ \mathrm{mol/L}$，在相同条件下测得吸光度为$0.587$。
    \end{enumerate}
    计算$\ce{MX3}$络合物的稳定常数$K_{\ce{MX3}}$。
\end{enumerate}